\begingroup\raggedright
\section*{A route through Chapter XI}
\phantomsection
\addcontentsline{toc}{section}{A route through Chapter XI}
This chapter presents the finite packet theory and retains the detailed restricted arithmetic constructions. Two separate research companions explain the later progress toward Conjecture 66. They do not replace this chapter's proved statements or turn its conditional comparison into an existence theorem.

\subsection*{The mathematical text and its companions}
\begin{tabular}{@{}>{\raggedright\arraybackslash}p{.28\linewidth}>{\raggedright\arraybackslash}p{.67\linewidth}@{}}\toprule
Reader & Route and purpose\\\midrule
Mathematical text & Finite tame and wild theory, restricted arithmetic constructions, and the exact realization/comparison boundary.\\[7pt]
Core visual guide & \href{https://komorebi.invalid/b07/core_visual}{Geometric route through the mathematical text}. The original arithmetic examples retain their stated scope.\\[7pt]
Formal companion & \href{https://komorebi.invalid/b07/core_lean}{Formal statements, literal excerpts and evidence lineage}. The guide distinguishes the selected source census, compiler replay and arithmetic hypotheses.\\[7pt]
Progress paper & \href{https://komorebi.invalid/b07/research_start}{Toward Arithmetic Packet Realization}. Connected proofs, failed routes, conductor conventions and conditions on native markings.\\[7pt]
Research visual atlas & \href{https://komorebi.invalid/b07/atlas_start}{Arithmetic Packet Realization}. Thirty-two worked plates accompany the progress paper.\\\bottomrule
\end{tabular}

\subsection*{The endpoint is a quantified construction problem}
Conjecture 66 asks for an actual arithmetic witness for each admitted source with its supplied H1--H4 datum. The comparison proposition then uses H5 for that same realized candidate. The formal companion distinguishes this conjecture from the differently numbered statements in the source registers.

The \href{https://komorebi.invalid/b07/contract}{research contract in Section 1} and \href{https://komorebi.invalid/b07/contract_atlas}{Plate XI.1} make the quantifiers explicit. Supplying finite operators does not construct an arithmetic representation. Likewise, proving a factor identity for a different candidate does not identify the packet's realizer.

\subsection*{How to use the links}
The reading links connect the mathematical text, visual studies, research companions and formal companion. Section and result numbers remain local to the named treatment. Plate references include the chapter number, so a reference such as Plate XI.32 identifies a visual plate rather than a source version.
\clearpage
\section*{Scope and research routes}
\phantomsection
\addcontentsline{toc}{section}{Scope and research routes}
\begin{tabular}{@{}>{\raggedright\arraybackslash}p{.32\linewidth}>{\raggedright\arraybackslash}p{.63\linewidth}@{}}\toprule
Question & Route beyond the core statement\\\midrule
Parameter, full character and recovery & \href{https://komorebi.invalid/b07/arithmetic}{Progress Sections 7--10} and \href{https://komorebi.invalid/b07/arithmetic_atlas}{Plates XI.6--XI.10}. The full character, central character and noncompact completion remain fixed.\\[8pt]
Monomial, Clifford and occurrence routes & \href{https://komorebi.invalid/b07/monomial}{Progress Part II}. The original character route is excluded under its hypotheses. A different corrected character passes a scalar test, but occurrence in the same candidate is still open.\\[8pt]
Conductor and capacity & \href{https://komorebi.invalid/b07/conductor}{Progress Sections 22--25} and \href{https://komorebi.invalid/b07/conductor_atlas}{Plates XI.11--XI.19}. Swan, ordinary conductor and epsilon exponent for an additive character are named separately. Equal numerical invariants need not give equal quotient capacity.\\[8pt]
Native markings and allowed repairs & \href{https://komorebi.invalid/b07/markings}{Progress Sections 26--29} and \href{https://komorebi.invalid/b07/markings_atlas}{Plates XI.20--XI.29}. A freely transported marking is different from a prescribed native marking. Finite corrections require their stated inputs and permission to change the relevant map.\\[8pt]
Remaining arithmetic certificate & \href{https://komorebi.invalid/b07/remaining}{Progress Section 30} and \href{https://komorebi.invalid/b07/remaining_atlas}{Plates XI.30--XI.32}. General realization, H5, corrected-character occurrence and the stronger native certificate remain separate tasks.\\\bottomrule
\end{tabular}

\subsection*{The construction and its limits}
The breakwise phase product uses its moving symplectic spaces. The type-side
pairing fixes a prime-field character, and the tame comparison retains the
boundaries for the rank-one conductor and the Gauss sum of the trivial character. Interaction
carriers are distinguished from minimal character orbits.

The detailed construction appendices give ordinary proofs under their stated
type-theory hypotheses. The selected Henniart--Lemaire induction inputs
\cite{HenniartLemaire2011} retain their characteristic, rank and normalization
conditions. They do not establish H5 or corrected-character occurrence.
The formal companion identifies the actual compiler and axiom coverage.

\endgroup
